\documentclass[11pt]{article}
\usepackage{amsmath,amssymb,amsthm,hyperref,geometry}\geometry{margin=1in}
\newtheorem{theorem}{Theorem}
\title{Fubini for directed and invertible grade contraction}
\date{15 November 2026}\begin{document}\maketitle
\paragraph{Status.} Independently reviewed general presentation theorem. This concerns contraction of retained category-valued grade diagrams. It is not yet a semantic theorem about arbitrary concurrent systems. The universal-colimit mechanism is established background; no novelty is asserted. Kengo Hirata, \emph{Notes on Lax Ends}, \url{https://arxiv.org/pdf/2210.01522}, Proposition 3.10, proves lax end/coend Fubini; Proposition 2.4 records the cross-variable compatibility that is essential here. The argument below spells out the retained generators and selected invertible-coordinate variant for this application.

For a strict functor $F:S^{op}\times S\to\mathbf{Cat}$ write $\mathsf L_SF$ for the directed-balance category: diagonal objects, native arrows and a directed slide for every state arrow, with naturality, identity and composition relations. Write $\mathsf P_SF$ for the same presentation with all slides inverted. The latter is the invertible-balance version, not an ordinary set-valued quotient. Universe sizes are fixed so all categories and presentations are small at the chosen level.

\begin{theorem}[Mixed Fubini]
Let $F:S^{op}\times S\times T^{op}\times T\to\mathbf{Cat}$ be strict. Choose independently whether $S$-slides and $T$-slides are directed or invertible. The two iterated contractions are canonically isomorphic to the simultaneous presentation with both types of slides and their interchange relation. With both directed this is $\mathsf L_{S\times T}\widetilde F$, where
\[
 \widetilde F((s,t),(s',t'))=F(s,s',t,t').
\]
If some types are invertible, localize exactly those coordinate slides. The comparisons preserve all native arrows and the chosen direction of every noninvertible slide.
\end{theorem}
\begin{proof}
An object of either iterated presentation is the same triple $[s,t,x]$ with $x\in F(s,s,t,t)$. Its generators are native arrows, $S$-slides and $T$-slides. Each inner presentation has its own defining relations. Outer slide naturality must hold on \emph{every} inner arrow, hence on inner slides as well. That is precisely the mixed interchange relation. Conversely imposing this relation makes the outer transformations natural on paths and on the inner quotient, so produces exactly the iterated presentation. This proves an isomorphism by inverse assignments on generators, not by merely comparing object sets.

For the product-index presentation, a slide for $(h,k)$ factors by the composition relation into slides for $(h,1)$ and $(1,k)$. Factoring in the other order gives their interchange relation because the two product-category composites agree. Thus coordinate slides generate the same presentation with the same relations. Native arrows are fixed throughout.

Inverting a selected family adds inverse generators and their inverse equations. The previous generator identifications preserve exactly that family. Therefore the same isomorphisms descend to the mixed directed/invertible versions. No inversion of the other coordinate slides is implied.
\end{proof}

\begin{theorem}[Finite contraction coherence]
For finitely many independently typed grade variables, every order and bracketing of their directed/invertible contractions gives the same simultaneous category up to the canonical comparison. All diagrams of these canonical comparison functors commute, including their action on native arrows.
\end{theorem}
\begin{proof}
Apply the preceding generator description to the finite tuple of all variables. The common presentation has native arrows, one slide family per coordinate, each coordinate's relations and pairwise interchange. Every canonical rebracketing or permutation comparison sends an addressed native generator or coordinate slide to that same generator in the simultaneous presentation. Hence any two composites of these comparisons agree as functors. This asserts coherence of the canonical comparisons only; it does not equate arbitrary arrows of the presented category.
\end{proof}

\section{Why ordinary horizontal units cannot be replaced indiscriminately}
The Fubini theorem does not say that directed contraction can replace every relative tensor in the original factory. Let $S$ be $0\to1$ and put $F(s,t)=S(0,t)$, as a discrete category, with the other variable constant. Both values are singletons, so the directed contraction is exactly $S$: its nonidentity balance remains $0\to1$. It is not equivalent to the terminal category. Ordinary coYoneda contraction of the same representable against the constant singleton gives the terminal category; invertible balance also makes this particular example equivalent to terminal.

Thus replacing the existing pure-middle pseudocoend everywhere by directed contraction does not preserve the usual representable unit up to equivalence. It would require a different, potentially skew, unital doctrine. The current proposal instead keeps pure horizontal balance invertible and marks only the additional shared-grade contractions as directed. This is a structural distinction needed to preserve the existing relative unifier, not merely a stylistic choice of notation. It is not by itself a complete repair for arbitrary network composition: if an added directed grade wire is later exposed as a compositional boundary, its proposed identity node faces the same coYoneda test. Such a wire needs a separate unital/evaluation construction or an explicitly different weak-unit doctrine. The actual directed state-extension construction instead composes using the existing process composition and tensor of retained state objects; its state unit is the monoidal unit, not an unjustified directed coend representable.

\section{The exact connection with shared wiring}
Independent context and content grades in the original comma need not remain independent in every extension. Linking an output grade to an input grade introduces a shared variable, and directed contraction retains its reparameterization maps. The preceding theorem permits several such links without selecting an order-dependent quotient procedure. It also keeps this operation distinct from existing pure-middle invertible balance.

A finite network can supply its diagram locally: choose a category-valued functor for each node, with independently typed input and output grade ports. For each internal wire pair one covariant port with one contravariant port of the same grade category. Form the product of the local node categories as a functor of all port variables, then contract the internal wire variables. Contractions of a subnetwork followed by contraction of the remaining wires agree with the simultaneous construction by mixed Fubini. Local witnesses and native arrows remain generators; this is not a proof that every inclusion is faithful, which needs a separate detector or other argument.

These are categories of presented network data. They are not automatically evaluated processes in the original $H$, nor probability laws determined by node marginals. Evaluation requires the relevant tensor, composition and feedback operations, with their compatibility. Multiway sharing requires explicit multiport node data or an admitted copying/joint-witness operation; pairing binary ports does not silently add it. Infinite networks require separately justified size/completion and evaluation conditions. These distinctions prevent an order-independent wiring presentation from being mistaken for a full unrestricted-concurrency theorem.

\section{Weak diagrams}
For pseudofunctorial input, slide composition includes its actual compositor/unit comparisons. The same universal-property argument suggests a canonical equivalence of weak presentations, rather than the literal strict isomorphisms above. To use it in the existing weak relative factory, one must write those comparisons and verify their typed interchange. That extension is an explicit next obligation, not an appeal to an unspecified strictification preserving all extra data.
\end{document}
